\documentclass[UTF8,zihao=-4]{ctexart}
\usepackage[a4paper,margin=2.2cm]{geometry}
\usepackage{fontspec}
\usepackage{tabularx}
\usepackage{array}
\usepackage{ragged2e}
\usepackage{fancyhdr}
\usepackage{hyperref}
\setCJKmainfont{Noto Serif SC}
\setCJKsansfont{Noto Sans SC}
\setmainfont{Noto Serif SC}
\setlength{\parindent}{0pt}
\setlength{\parskip}{0.45em}
\pagestyle{fancy}
\fancyhf{}
\fancyhead[L]{商务智能}
\fancyfoot[C]{\thepage}
\hypersetup{colorlinks=true,linkcolor=black,urlcolor=blue}
\begin{document}
\title{19、数据挖掘的方法}
\author{}
\date{}
\maketitle

定义：对数据库或数据仓库中蕴涵的、未知的、非平凡的、有潜在应用价值的模式或规则的提取\par

\begin{tabularx}{\linewidth}{|*{2}{>{\RaggedRight\arraybackslash}X|}}
\hline
方法 & 作用 \\
\hline
特征规则挖掘 & 描述一类数据对象的普遍特征 \\
\hline
关联规则挖掘 & 发现事务数据中项与项之间的关联关系 \\
\hline
序列模式分析 & 发现数据对象之间的前后时序关系 \\
\hline
分类分析 & 根据训练样本建立模型，对未知数据分类 \\
\hline
聚类分析 & 对没有类别标记的数据自动划分群组 \\
\hline
\end{tabularx}

1. 特征规则挖掘\par
是一种常见的知识形式，它用于描述一类数据对象的普遍特征，是普化知识的一种\par

特征规则的数据挖掘方法有两类：\par

1. 面向属性归约方法\par

它通过属性值之间的概念层次结构进行归约，获得概括性知识。例如，把“南京”提升为“江苏”，再把“江苏”提升为“华东”，这叫概念提升。概念提升的目的包括规范化属性取值，提高模式的置信度和兴趣度。\par

2. 数据立方方法\par

数据立方方法利用预先计算好的统计、求和等综合结果进行挖掘，可以减少重复统计查询，提高挖掘效率。它和 OLAP 中的上卷、下钻思想有关。\par

特征规则挖掘用于描述一类对象的普遍特征，属于普化知识发现。其主要方法包括面向属性归约方法和数据立方方法。面向属性归约通过概念层次树对属性值进行概念提升，以获得概括性知识。\par

2. 关联规则挖掘\par

关联规则挖掘用于发现事务数据库中多个属性或项之间的关联程度。\par

典型形式是：\par

A → B\par

含义是：出现 A 的事务中，经常也出现 B。\par

例如：\par

购买尿布 → 购买啤酒\par
购买面包和黄油 → 购买牛奶\par

关联规则挖掘用于发现事务数据中项与项之间的关联关系，规则形式为 A→B。其核心指标是支持度和置信度，满足最小支持度和最小置信度的规则称为强规则。Apriori 算法利用“频繁项集的任一子集必定频繁”的性质寻找频繁项集。\par

3. 序列模式分析\par

序列模式分析与关联规则类似，也是为了找出数据对象之间的联系，但它更强调前后时序关系。也就是说，它不仅关心“哪些项经常一起出现”，还关心“哪些事件先发生、哪些事件后发生”。\par

例如：\par

下雨 → 洪涝\par
购买电筒 → 购买电池\par

序列模式分析用于发现具有时间先后关系的数据对象之间的联系。它与关联规则挖掘类似，但更强调事件之间的前因后果和时序关系。\par

4. 分类分析\par

分类分析是数据挖掘的重要方法之一。分类分析通过分析训练数据样本，产生关于类别的精确描述，并用于对未来数据进行分类和预测。\par

1. 基本思想\par

分类分析需要先有已经标记类别的数据。例如：\par

信用等级：优、良、一般、差；\par
是否购买电脑：yes / no；\par
客户是否流失：流失 / 未流失。\par

系统先根据这些带类别标记的数据学习规则，再对未知类别的数据进行预测。\par

2. 分类分析的两个步骤\par
\begin{tabularx}{\linewidth}{|*{2}{>{\RaggedRight\arraybackslash}X|}}
\hline
步骤 & 内容 \\
\hline
建立模型 & 用训练数据集构造分类模型 \\
\hline
使用模型分类 & 评估模型准确性，并对未知类别数据分类 \\
\hline
\end{tabularx}

因为训练样本已经给出了类标号，所以分类分析属于有指导学习。\par

3. 常见分类方法\par

包括：\par
决策树方法\par
贝叶斯方法\par
人工神经网络方法\par
约略集方法\par
遗传算法\par

其中，决策树是一种树结构：内部结点表示属性测试，边表示测试结果，叶结点表示类别或类别分布。\par

分类分析是一种有指导学习方法。它先利用带有类标号的训练数据建立分类模型，再用该模型对未知类别的数据进行分类和预测。常见方法包括决策树、贝叶斯、人工神经网络、约略集和遗传算法。\par

5. 聚类分析\par

聚类分析又称集群分析。它和分类分析相反：分类分析的数据已经有类别标记，而聚类分析面对的是没有类别标记的记录。系统按照距离或相似性等规则自动划分记录集合，相当于自动给数据打标签。\par

聚类分析可以分为距离聚类和相似系数聚类，主要方法包括：\par

划分方法\par
层次方法\par
基于密度的方法\par
基于网格的方法\par
基于模型的方法\par

聚类分析的结果可以表现为聚类树\par

聚类分析是一种无指导学习方法，输入数据没有预先给定类别标记，系统根据距离或相似性自动把记录划分为若干簇。主要方法包括划分方法、层次方法、基于密度的方法、基于网格的方法和基于模型的方法。\par

总结：常用的数据挖掘方法包括特征规则挖掘、关联规则挖掘、序列模式分析、分类分析和聚类分析等。特征规则挖掘用于描述一类数据对象的普遍特征，常用方法有面向属性归约方法和数据立方方法；关联规则挖掘用于发现事务数据中项与项之间的关联关系，规则形式为 A→B，核心指标是支持度和置信度；序列模式分析用于发现数据对象之间的前后时序关系；分类分析是一种有指导学习方法，通过带类标号的训练数据建立模型，并对未知类别数据进行分类和预测；聚类分析是一种无指导学习方法，对没有类别标记的数据按照距离或相似性自动划分为若干簇。此外，数据挖掘还包括异常分析、时序数据挖掘等方法。\par

\end{document}
